\documentclass{article}
\usepackage{amsmath,amssymb}
\usepackage{xurl}
\usepackage[hidelinks]{hyperref}
\setlength{\emergencystretch}{2em}
% Shared commands for the notes in docs/paper-gaps/.
% Every note loads this file via % Shared commands for the notes in docs/paper-gaps/.
% Every note loads this file via % Shared commands for the notes in docs/paper-gaps/.
% Every note loads this file via \input{command}. Project-wide notation
% belongs here, never in per-note redefinitions.

% --- Project configuration ------------------------------------------------
\newcommand{\projectIssueBase}{https://github.com/LionSR/QICLean/issues/}
\newcommand{\projectPrBase}{https://github.com/LionSR/QICLean/pull/}
\newcommand{\projectDocBase}{https://sirui-lu.com/QICLean/docs/}
% ---------------------------------------------------------------------------

\DeclareUnicodeCharacter{2080}{\ifmmode{}_{0}\else\textsubscript{0}\fi}
\DeclareUnicodeCharacter{2081}{\ifmmode{}_{1}\else\textsubscript{1}\fi}
\DeclareUnicodeCharacter{2082}{\ifmmode{}_{2}\else\textsubscript{2}\fi}
\DeclareUnicodeCharacter{2099}{\ifmmode{}_{n}\else\textsubscript{n}\fi}

\newcommand{\Lean}{\textsc{Lean}}

% A formal declaration name, upright sans, breakable at underscores.
\ExplSyntaxOn
\NewDocumentCommand{\leanid}{m}{
  \ifmmode
    \textsf{\detokenize{#1}}
  \else
    \group_begin:
    \urlstyle{sf}
    \tl_set:Nn \l_tmpa_tl {#1}
    \tl_replace_all:Nnn \l_tmpa_tl {\_} {_}
    \exp_args:NV \path \l_tmpa_tl
    \group_end:
  \fi
}
\ExplSyntaxOff

% Traceability links; use in footnotes, not in the running prose.
\newcommand{\ghissue}[1]{\href{\projectIssueBase#1}{GitHub issue \##1}}
\newcommand{\ghpr}[1]{\href{\projectPrBase#1}{PR \##1}}
\newcommand{\doclink}[1]{\href{\projectDocBase#1.html}{\path{#1}}}

% Domain notation, using \providecommand so a note-local refinement stays
% possible.
\providecommand{\tr}{\operatorname{tr}}

% Dirac notation and the identity operator, as in blueprint/src/macros.
\usepackage{dsfont}
\providecommand{\ket}[1]{|#1\rangle}
\providecommand{\bra}[1]{\langle #1|}
\providecommand{\braket}[2]{\langle #1 | #2 \rangle}
\providecommand{\ketbra}[2]{|#1\rangle\!\langle #2|}
\providecommand{\Id}{\mathds{1}}
\providecommand{\one}{\mathds{1}}
\providecommand{\C}{\mathbb{C}}
\providecommand{\MN}[1]{M_{#1}(\C)}
\providecommand{\MD}{M_D(\C)}

% Verdict marker: \gapnote{<kind>}{<status>}. Kind is the policy
% classification (clarification, local-correction, scope-restriction,
% unfaithful, false-source, open-gap); status is open, wip (elimination
% underway), resolved, or historical. Machine-read by the site index;
% prints a verdict line.
\newcommand{\gapnote}[2]{\par\noindent\textbf{Verdict:} #1 (#2).\par}
. Project-wide notation
% belongs here, never in per-note redefinitions.

% --- Project configuration ------------------------------------------------
\newcommand{\projectIssueBase}{https://github.com/LionSR/QICLean/issues/}
\newcommand{\projectPrBase}{https://github.com/LionSR/QICLean/pull/}
\newcommand{\projectDocBase}{https://sirui-lu.com/QICLean/docs/}
% ---------------------------------------------------------------------------

\DeclareUnicodeCharacter{2080}{\ifmmode{}_{0}\else\textsubscript{0}\fi}
\DeclareUnicodeCharacter{2081}{\ifmmode{}_{1}\else\textsubscript{1}\fi}
\DeclareUnicodeCharacter{2082}{\ifmmode{}_{2}\else\textsubscript{2}\fi}
\DeclareUnicodeCharacter{2099}{\ifmmode{}_{n}\else\textsubscript{n}\fi}

\newcommand{\Lean}{\textsc{Lean}}

% A formal declaration name, upright sans, breakable at underscores.
\ExplSyntaxOn
\NewDocumentCommand{\leanid}{m}{
  \ifmmode
    \textsf{\detokenize{#1}}
  \else
    \group_begin:
    \urlstyle{sf}
    \tl_set:Nn \l_tmpa_tl {#1}
    \tl_replace_all:Nnn \l_tmpa_tl {\_} {_}
    \exp_args:NV \path \l_tmpa_tl
    \group_end:
  \fi
}
\ExplSyntaxOff

% Traceability links; use in footnotes, not in the running prose.
\newcommand{\ghissue}[1]{\href{\projectIssueBase#1}{GitHub issue \##1}}
\newcommand{\ghpr}[1]{\href{\projectPrBase#1}{PR \##1}}
\newcommand{\doclink}[1]{\href{\projectDocBase#1.html}{\path{#1}}}

% Domain notation, using \providecommand so a note-local refinement stays
% possible.
\providecommand{\tr}{\operatorname{tr}}

% Dirac notation and the identity operator, as in blueprint/src/macros.
\usepackage{dsfont}
\providecommand{\ket}[1]{|#1\rangle}
\providecommand{\bra}[1]{\langle #1|}
\providecommand{\braket}[2]{\langle #1 | #2 \rangle}
\providecommand{\ketbra}[2]{|#1\rangle\!\langle #2|}
\providecommand{\Id}{\mathds{1}}
\providecommand{\one}{\mathds{1}}
\providecommand{\C}{\mathbb{C}}
\providecommand{\MN}[1]{M_{#1}(\C)}
\providecommand{\MD}{M_D(\C)}

% Verdict marker: \gapnote{<kind>}{<status>}. Kind is the policy
% classification (clarification, local-correction, scope-restriction,
% unfaithful, false-source, open-gap); status is open, wip (elimination
% underway), resolved, or historical. Machine-read by the site index;
% prints a verdict line.
\newcommand{\gapnote}[2]{\par\noindent\textbf{Verdict:} #1 (#2).\par}
. Project-wide notation
% belongs here, never in per-note redefinitions.

% --- Project configuration ------------------------------------------------
\newcommand{\projectIssueBase}{https://github.com/LionSR/QICLean/issues/}
\newcommand{\projectPrBase}{https://github.com/LionSR/QICLean/pull/}
\newcommand{\projectDocBase}{https://sirui-lu.com/QICLean/docs/}
% ---------------------------------------------------------------------------

\DeclareUnicodeCharacter{2080}{\ifmmode{}_{0}\else\textsubscript{0}\fi}
\DeclareUnicodeCharacter{2081}{\ifmmode{}_{1}\else\textsubscript{1}\fi}
\DeclareUnicodeCharacter{2082}{\ifmmode{}_{2}\else\textsubscript{2}\fi}
\DeclareUnicodeCharacter{2099}{\ifmmode{}_{n}\else\textsubscript{n}\fi}

\newcommand{\Lean}{\textsc{Lean}}

% A formal declaration name, upright sans, breakable at underscores.
\ExplSyntaxOn
\NewDocumentCommand{\leanid}{m}{
  \ifmmode
    \textsf{\detokenize{#1}}
  \else
    \group_begin:
    \urlstyle{sf}
    \tl_set:Nn \l_tmpa_tl {#1}
    \tl_replace_all:Nnn \l_tmpa_tl {\_} {_}
    \exp_args:NV \path \l_tmpa_tl
    \group_end:
  \fi
}
\ExplSyntaxOff

% Traceability links; use in footnotes, not in the running prose.
\newcommand{\ghissue}[1]{\href{\projectIssueBase#1}{GitHub issue \##1}}
\newcommand{\ghpr}[1]{\href{\projectPrBase#1}{PR \##1}}
\newcommand{\doclink}[1]{\href{\projectDocBase#1.html}{\path{#1}}}

% Domain notation, using \providecommand so a note-local refinement stays
% possible.
\providecommand{\tr}{\operatorname{tr}}

% Dirac notation and the identity operator, as in blueprint/src/macros.
\usepackage{dsfont}
\providecommand{\ket}[1]{|#1\rangle}
\providecommand{\bra}[1]{\langle #1|}
\providecommand{\braket}[2]{\langle #1 | #2 \rangle}
\providecommand{\ketbra}[2]{|#1\rangle\!\langle #2|}
\providecommand{\Id}{\mathds{1}}
\providecommand{\one}{\mathds{1}}
\providecommand{\C}{\mathbb{C}}
\providecommand{\MN}[1]{M_{#1}(\C)}
\providecommand{\MD}{M_D(\C)}

% Verdict marker: \gapnote{<kind>}{<status>}. Kind is the policy
% classification (clarification, local-correction, scope-restriction,
% unfaithful, false-source, open-gap); status is open, wip (elimination
% underway), resolved, or historical. Machine-read by the site index;
% prints a verdict line.
\newcommand{\gapnote}[2]{\par\noindent\textbf{Verdict:} #1 (#2).\par}

\title{Boundary eigenvalues in the error reduction by a spectral projector}
\author{}
\date{2026-10-25}
\begin{document}
\maketitle
\gapnote{false-source}{resolved}

\section{The printed statement}

Let $\mathcal H=\mathcal H_L\otimes\mathcal H_M\otimes\mathcal H_R$ and let $H$
be a Hamiltonian on $\mathcal H$ with ground energy $\varepsilon_0$ and no
eigenvalue in the open interval $(\eta_0,\eta_1)$. Definition
\texttt{def:spectral-AGSP} of Arad--Landau--Vazirani--Vidick,
\href{https://arxiv.org/abs/1602.08828}{arXiv:1602.08828v2}, source lines
570--581, calls a positive semidefinite operator
$K=\sum_{i=1}^{D^2}L_i\otimes A_i\otimes R_i$ a $(D,\Delta)$-spectral
projector for $(H,\eta_0,\eta_1)$ if $H$ and $K$ have the same eigenvectors,
eigenvalues of $H$ smaller than $\eta_0$ correspond to eigenvalues of $K$ at
least $1$, and eigenvalues of $H$ larger than $\eta_1$ correspond to
eigenvalues of $K$ smaller than $\sqrt\Delta$. Lemma \texttt{l:er}, lines
597--622, asserts that if $S\subseteq\mathcal H_M$ is $\delta$-viable for
$T=H_{[\varepsilon_0,\eta_0]}$, then
$V=\operatorname{span}\{A_iS\}$ has dimension at most $D^2\dim S$ and is
$\delta'$-viable for $T$ with
\begin{align}
    \delta'&=\frac{\Delta}{(1-\delta)^2}.
    \label{eq:alvv_er_delta}
\end{align}

\section{A counterexample at the boundary}

The definition places no condition on an eigenvalue of $H$ equal to
$\eta_0$, but such an eigenvector lies in $T$.
Take $\mathcal H_L=\mathcal H_R=\mathbb C$, $\mathcal H_M=\mathbb C^2$,
$H=\operatorname{diag}(\eta_0,\eta_1+1)$, $D=1$ and $K=A_1=0$, with
$\Delta=1/4$. Every condition of the definition holds, since neither
eigenvalue of $H$ is smaller than $\eta_0$ and $0<\sqrt\Delta$. The space
$S=\mathbb C^2$ is $0$-viable for $T=\operatorname{span}\{e_0\}$, but
$V=\{0\}$ is not $\delta'$-viable for $T$ for any $\delta'<1$, whereas
(\ref{eq:alvv_er_delta}) gives $\delta'=1/4$. The same issue arises at
$\eta_1$: an eigenvector with eigenvalue exactly $\eta_1$ lies outside $T$
and carries no bound on its $K$-eigenvalue.

Definition \texttt{d:vs}, line 434, restricts viability to
$0\le\delta\le1$. At $\delta=1$ the parameter (\ref{eq:alvv_er_delta}) is
undefined, and the proof uses $\delta<1$ through Lemma \texttt{l:2}; this is
a degenerate reading rather than an error, and the formal statement assumes
$\delta<1$.

\section{The corrected statement}

The counterexample disappears when the spectral conditions are closed:
eigenvalues $h\le\eta_0$ of $H$ correspond to eigenvalues $k\ge1$ of $K$, eigenvalues $h\ge\eta_1$ correspond to
eigenvalues with $k^2\le\Delta$. With $\delta<1$, the proof of the source
then goes through. Positivity of $K$ is not used; dropping it matters
because the unsquared Chebyshev filter that the source later identifies with
such an operator is not positive semidefinite (see the companion note
\path{alvv16_chebyshev_filter_positivity.tex}).

The proof displays the bound
$\alpha'^2/(\alpha'^2+(1-\alpha'^2)\Delta)$ for the overlap of $Ku/\|Ku\|$
with $v$. Writing $u=\alpha v'+\sqrt{1-\alpha^2}\,v^\perp$, the squared norm
of $Ku$ is at most $\alpha'^2+(1-\alpha^2)\Delta$, so the denominator should
carry $1-\alpha^2$. The final bound $1-\Delta/(1-\delta)^2$ is unaffected.

The formal statement uses a common orthonormal eigenbasis $b_j$ of $H$ and
$K$, with eigenvalues $h_j$ and $k_j$, and takes
$T=\operatorname{span}\{b_j:h_j\le\eta_0\}$. It separates the operator step:
if $K$ maps $T$ and $T^\perp$ into themselves, $\|Kt\|\ge\|t\|$ on $T$ and
$\|Kq\|^2\le\Delta\|q\|^2$ on $T^\perp$, then $K(S\otimes\text{rest})$ is
$\delta'$-close to $T$. The corrected lemma is
\leanid{Submodule.IsViable.spectralProjector}, with the definition
\leanid{Submodule.IsSpectralProjector} and the operator step
\leanid{Submodule.IsCloseTo.map_of_contract}.
\end{document}
